\scanpage{38}
D'autre part, on aura ici
\[
\beta\in\unsure{Q|\mathbb Q}(\mathbb R)\subset P'(\mathbb R)
\]
donc
\[
\mathbb R(\gamma)=\mathbb R(\beta,\gamma)=\mathbb R(\beta,\alpha)=\mathbb R(\alpha)
\]
et par suite
\[
\mathbb R(\gamma)=\mathbb R(\alpha)
\]
qui est donc $=\mathbb C$ sauf dans le cas exceptionnel précisé
plus haut.

\uline{Remarque} Soit $M$ un motif sur $\mathbb R$, correspondant à un \deleted{\unsure{$G$-module}
\unsure{aussi noté}} $T_B(M)$. Alors on ne peut avoir $f_{\infty M}=\mathrm{id}_M$ que
si $M$ est une somme de puissances tensorielles
de $\mathbb Q(2)$.

\deleted{La relation
\[
 f_\infty j(\lambda) f_\infty^{-1}=j(\lambda')
\]
implique, faisant $\lambda=C\in\unsure{S|D}(\mathbb R$
\[
f_\infty C f_\infty^{-1}=C^{-1}=\unsure{\eta|y}C\qquad(\eta=i(-1))
\]
donc
\[
 f_\infty\in\underline{\mathrm{Norm}}_{G_{\mathbb R}}(C,C^{-1})
\]
\deleted{\unsure{Or, on a} $K'=\underline{\mathrm{Norm}}_{G_{\mathbb R}}(C,C^{-1})\supset\underline{\mathrm{Norm}}_{G_{\mathbb R}}(C)=K$ qui est un
sous groupe d'indice 1 ou 2. On sait que $K(\mathbb R)$ est un
compact maximal dans $G(\mathbb R)$} Considérons $G/i(G_m)=\dot G$
$[=\mathrm{Ker}\,\varepsilon/i(\mu_2)]$, et les images \deleted{\unsure{de}} $\dot f_\infty$, $\dot C$ de $f_\infty$, $C$,
donc on a
\[
 \dot f_\infty\in\underline{\mathrm{Norm}}_{\dot G_{\mathbb R}}(\dot C,\dot C^{-1})\overset{\mathrm{dfn}}=K'
\]
Or \deleted{on sait \unsure{qu'on}} sait
\[
 K=\underline{\mathrm{Norm}}_{\dot G_{\mathbb R}}(\dot C)\subset K'
\]}
% archive sheet number: 37
